\chapter{Pineal Gland}
\index{Pineal Gland}

\medskip
\noindent\textit{This chapter is for general education. A scan finding near the pineal gland should be interpreted by the clinician who ordered the scan.}

\section*{What it is}
The pineal gland is a small structure deep in the brain. It produces melatonin, a hormone involved in the timing of sleep and wakefulness. Its activity is influenced by the light--dark cycle.

\section*{Pineal-region findings}
Small cysts near the pineal gland are often found incidentally on brain imaging and may cause no symptoms. Less commonly, a mass or another problem in this region can affect nearby structures or the flow of cerebrospinal fluid. Symptoms depend on the size and nature of the finding, not on the name alone.

\section*{Assessment}
Clinicians interpret the scan together with symptoms and examination. Follow-up MRI, eye or neurological assessment, hormone tests, or specialist referral may be advised. Do not assume that a pineal cyst is the cause of headaches or sleep problems without medical review.

\section*{When to Seek Medical Care}
Seek emergency help for sudden severe headache, repeated vomiting, new weakness, confusion, seizure, double vision, loss of consciousness, or rapidly worsening balance. Discuss persistent headaches, visual changes, or a scan finding with the treating clinician.

\section*{Sources and further reading}
\begin{itemize}
\item National Institute of Neurological Disorders and Stroke. \textit{Brain basics: understanding sleep}. \url{https://www.ninds.nih.gov/health-information/public-education/brain-basics/brain-basics-understanding-sleep}
\item Radiopaedia. \textit{Pineal gland cyst}. \url{https://radiopaedia.org/articles/pineal-gland-cyst}
\end{itemize}
